%%=============================================================================
%% Inleiding
%%=============================================================================

\chapter{\IfLanguageName{dutch}{Inleiding}{Introduction}}%
\label{ch:inleiding}

In de hedendaagse digitale samenleving zijn digitale bestanden alomtegenwoordig, zowel in privé- als
in bedrijfscontext. Documenten, onderzoeksgegevens en administratieve dossiers worden digitaal
opgemaakt, bijgehouden en wereldwijd verspreid. Deze digitalisering vergemakkelijkt het beheer van
informatie, maar brengt ook een kritieke afhankelijkheid met zich mee: toegang tot die bestanden
vereist een betrouwbare netwerkverbinding.

Die verbinding is echter niet altijd gegarandeerd. In afgelegen gebieden, op zee of tijdens
rampen waarbij de infrastructuur wegvalt, is conventionele internetconnectiviteit onvoldoende of
volledig afwezig~\autocite{Xia2022}. Bestaande alternatieven zoals commerciële satellietdiensten
zijn voor particulieren en kleine organisaties vaak financieel of praktisch ontoereikend.

Radioverbindingen bieden een potentieel alternatief. Ze zijn infrastructuuronafhankelijk, hebben
een breed gedocumenteerde geschiedenis in noodcommunicatie, en vereisen relatief eenvoudige
hardware. Toch is de toepasbaarheid van radioverbindingen voor toegang tot een persoonlijke
fileserver nog niet grondig onderzocht. Deze bachelorproef vult die lacune op.

\section{\IfLanguageName{dutch}{Probleemstelling}{Problem Statement}}%
\label{sec:probleemstelling}

% Uit je probleemstelling moet duidelijk zijn dat je onderzoek een meerwaarde heeft voor een concrete doelgroep. De doelgroep moet goed gedefinieerd en afgelijnd zijn. Doelgroepen als ``bedrijven,'' ``KMO's'', systeembeheerders, enz.~zijn nog te vaag. Als je een lijstje kan maken van de personen/organisaties die een meerwaarde zullen vinden in deze bachelorproef (dit is eigenlijk je steekproefkader), dan is dat een indicatie dat de doelgroep goed gedefinieerd is. Dit kan een enkel bedrijf zijn of zelfs één persoon (je co-promotor/opdrachtgever).

De doelgroep van dit onderzoek bestaat uit IT-professionals en netwerkingenieurs die actief zijn
in omgevingen met beperkte of onbetrouwbare internetconnectiviteit. Concreet gaat het om:

\begin{itemize}
  \item maritieme operatoren die op zee werken en geen stabiele breedbandverbinding hebben;
  \item hulpverleningsorganisaties die in rampgebieden opereren waar de netwerkinfrastructuur
        is uitgevallen;
  \item bedrijven of onderzoekers die actief zijn in afgelegen regio's zonder
        vaste netwerkinfrastructuur;
  \item radioamateurs die bestaande radioinfrastructuur willen uitbreiden met
        netwerktoepassingen.
\end{itemize}

Voor deze groepen zijn bestaande oplossingen zoals Starlink of vaste breedbandverbindingen
geen haalbare optie. Er is behoefte aan een methode om via een zelfopgezette radioverbinding
toegang te krijgen tot een eigen fileserver, zonder afhankelijkheid van commerciële of
gecentraliseerde infrastructuur.

\section{\IfLanguageName{dutch}{Onderzoeksvraag}{Research question}}%
\label{sec:onderzoeksvraag}

% Wees zo concreet mogelijk bij het formuleren van je onderzoeksvraag. Een onderzoeksvraag is trouwens iets waar nog niemand op dit moment een antwoord heeft (voor zover je kan nagaan). Het opzoeken van bestaande informatie (bv. ``welke tools bestaan er voor deze toepassing?'') is dus geen onderzoeksvraag. Je kan de onderzoeksvraag verder specifiëren in deelvragen. Bv.~als je onderzoek gaat over performantiemetingen, dan 

De centrale onderzoeksvraag luidt:

\begin{quote}
  \textit{Hoe kan een verbinding worden opgezet met een eigen fileserver via een
  radioverbinding, en wat zijn de technische randvoorwaarden en maximale
  afstandsbeperkingen voor een werkende implementatie?}
\end{quote}

Deze vraag wordt verder opgesplitst in de volgende deelvragen:

\begin{itemize}
  \item Welke radioprotocollen (zoals AX.25, Winlink of IP-over-radio) zijn geschikt voor
        bestandsoverdracht, en wat zijn hun kenmerken op het vlak van bandbreedte,
        betrouwbaarheid en bereik?
  \item Welke multiplextechnieken zijn van toepassing binnen radiogebaseerde
        datacommunicatie, en welke beperkingen leggen zij op?
  \item Wat is de maximale praktische afstand waarover een bestandsoverdracht via
        radioverbinding haalbaar is, afhankelijk van frequentieband en omgevingsfactoren?
  \item Welke minimale infrastructuur en configuratie zijn vereist voor een werkend
        proof-of-concept?
\end{itemize}

\section{\IfLanguageName{dutch}{Onderzoeksdoelstelling}{Research objective}}%
\label{sec:onderzoeksdoelstelling}

% Wat is het beoogde resultaat van je bachelorproef? Wat zijn de criteria voor succes? Beschrijf die zo concreet mogelijk. Gaat het bv.\ om een proof-of-concept, een prototype, een verslag met aanbevelingen, een vergelijkende studie, enz.

Het beoogde eindresultaat van deze bachelorproef is een werkende proof-of-concept waarbij een
fileserver bereikbaar wordt gemaakt via een punt-tot-punt radioverbinding. Het succes van het
onderzoek wordt beoordeeld aan de hand van meetbare criteria: de gerealiseerde bandbreedte,
het packet loss percentage, de verbindingsstabiliteit en de maximale operationele afstand.

Indien een volledig werkende opstelling niet haalbaar blijkt, wordt dit beschouwd als een
waardevol resultaat op zich: de bachelorproef levert dan een theoretisch onderbouwde verklaring
op van de technische grenzen van radiogebaseerde bestandsoverdracht, aangevuld met
aanbevelingen voor vervolgonderzoek of alternatieve benaderingen.

\section{\IfLanguageName{dutch}{Opzet van deze bachelorproef}{Structure of this bachelor thesis}}%
\label{sec:opzet-bachelorproef}

% Het is gebruikelijk aan het einde van de inleiding een overzicht te
% geven van de opbouw van de rest van de tekst. Deze sectie bevat al een aanzet
% die je kan aanvullen/aanpassen in functie van je eigen tekst.

De rest van deze bachelorproef is als volgt opgebouwd:

In Hoofdstuk~\ref{ch:stand-van-zaken} wordt een overzicht gegeven van de stand van zaken
binnen het onderzoeksdomein, op basis van een literatuurstudie. Aan bod komen radioprotocollen,
frequentiebanden, multiplextechnieken en bestaande systemen voor radiogebaseerde dataoverdracht.

In Hoofdstuk~\ref{ch:methodologie} wordt de methodologie toegelicht en worden de gebruikte
onderzoekstechnieken besproken: de requirementsanalyse, de opbouw van de proefopstelling
en de meetmethoden die gehanteerd worden om de onderzoeksvragen te beantwoorden.

% TODO: voeg hier hoofdstukken toe voor je proof-of-concept en resultaten

In Hoofdstuk~\ref{ch:conclusie}, tenslotte, wordt de conclusie gegeven en een antwoord
geformuleerd op de onderzoeksvragen. Daarbij wordt ook een aanzet gegeven voor toekomstig
onderzoek binnen dit domein.